\begin{center}
\textsc{Gauss} to \textsc{Bessel}.\quad Göttingen, 12~March 1826.
\end{center}

\par\noindent\hspace*{\parindent}.\ .\ .\ .\ .\ .\quad I spent most of last summer on my measurements in the Bremen and Oldenburg regions.\quad.\ .\ .\ .\ .\ .

\par\noindent\hspace*{\parindent}As for my measurements, I have completed the trigonometrical part, at least in the literal sense. My side Varel-Jever joins \textsc{Krayenhoff}'s triangles.
Whether they are really finished, however, I do not yet know myself.
My angle measurements at Jever show \textls[120]{quite enormous} differences from \textsc{Krayenhoff}'s, reaching \(15^{\prime\prime}\).
Their source
\par\clearpage
\noindent{\small\textbf{ON THE HANOVER TRIANGULATION.}}\hfill361\par\medskip
\noindent I cannot specify with certainty.
I vouch for my own angles to within an arcsecond.
The differences could be explained if I were entitled to assume that \textsc{Krayenhoff}'s center of the station lies one meter from mine (I observed one floor lower).
The information I received from \textsc{Krayenhoff} is unsatisfactory as regards centering, but in general it shows that his angle measurements in that region come nowhere near mine in precision.
He used an inferior instrument in East Frisia. Of many series of angles he always kept only those that seemed to fit best (without specifying by how much the others deviated), and even among those retained at Jever there are differences of \(4^{\prime\prime}\).
Under these circumstances I do not think it advisable to extend my measurements farther in East Frisia.
His measurements farther south are better, and a connection of my triangles through the Osnabrück region to Bentheim could undoubtedly give more reliable results.
But the uncertain state of my health has so far discouraged me from proposing operations of this kind, which would still require a campaign of a year or a year and a half, and scarcely anything substantial can be done in this regard this summer.\quad.\ .\ .\ .\ .\ .

